% ============================================================
% FOND COLORÉ SIDEBAR
% ============================================================
% Collé à la minipage (pas de ligne vide) : sinon ce tikzpicture forme un
% paragraphe à part et consomme une ligne complète (~12pt) en haut de page.
\begin{tikzpicture}[remember picture,overlay]
\fill[sidebarcolor] (current page.north west) rectangle ([xshift=4.9cm]current page.south west);
\end{tikzpicture}%
% Hauteur fixée avec position interne « s » (stretch) : les \vfill plus bas
% peuvent alors se détendre, ce qui centre l'icône Omarchy dans l'espace libre
% restant sous le contenu — sans overlay ni seconde passe de compilation.
% La hauteur disponible est \textheight moins \topskip : « top=0cm » dans la
% géométrie, c'est \topskip qui fait office de marge haute et décale d'autant
% la première ligne (donc ces deux minipages) vers le bas.
\begin{minipage}[t][\dimexpr\textheight-\topskip\relax][s]{0.20\textwidth}
    \vspace{0pt}
    % Texte en drapeau : la colonne est trop étroite pour une justification
    % propre (mots étirés, césures disgracieuses).
    \RaggedRight

    \begin{center}
\tikz\node[rounded corners=6pt, minimum width=3.6cm, minimum height=4.5cm, path picture={
    \node[yshift=-2mm] at (path picture bounding box.center){
        \includegraphics[width=3.6cm]{\myphoto}
    };
}, draw=maincolor, line width=2pt]{};
\end{center}

    \vspace{2pt}

    \subsection*{\textcolor{maincolor}{\faIcon{address-card}} Informations}
\small
\begin{tabular}{@{}l@{ }l@{}}
\hspace{8pt}\textcolor{maincolor}{\faIcon{map-marker-alt}} & \myaddress \\[3pt]
\hspace{8pt}\textcolor{maincolor}{\faIcon{phone}} & \myphone \\[3pt]
\hspace{8pt}\textcolor{maincolor}{\faIcon{envelope}} & \href{mailto:\myemail}{\scriptsize \myemail} \\[3pt]
\hspace{8pt}\textcolor{maincolor}{\faIcon{linkedin}} & \href{\mylinkedinurl}{\scriptsize \mylinkedinname} \\[3pt]
\hspace{8pt}\textcolor{maincolor}{\faIcon{github}} & \href{\mygithuburl}{\scriptsize \mygithubname}
\end{tabular}


    \vspace{10pt}

    \subsection*{\textcolor{maincolor}{\faIcon{tools}} Compétences}
\small

\hspace{4pt}\textbf{\color{maincolor}Back-end}
\begin{itemize}[leftmargin=11pt, itemsep=1pt, label={}]
    \item PHP / Symfony
    \item Go \textit{\scriptsize(perso)}
\end{itemize}

\vspace{2pt}
\hspace{4pt}\textbf{\color{maincolor}Front-end}
\begin{itemize}[leftmargin=11pt, itemsep=1pt, label={}]
    \item React / Vue
    \item React Native
\end{itemize}

\vspace{2pt}
\hspace{4pt}\textbf{\color{maincolor}Misc}
\begin{itemize}[leftmargin=11pt, itemsep=1pt, label={}]
    \item Docker
    \item Git
    \item Linux
\end{itemize}


    \vspace{6pt}

    \subsection*{\textcolor{maincolor}{\faIcon{star}} Qualités}
\small
\begin{itemize}[leftmargin=2pt, itemsep=1pt, label={}]
    \item Reproduire avant de corriger
    \item Lisibilité avant astuce
    \item Refactorisation
    \item Vulgarisation auprès des métiers
\end{itemize}


    \vspace{6pt}

    \subsection*{\textcolor{maincolor}{\faIcon{language}} Langues}
\small
\hspace{2pt}\textbf{Français} - Natif

\vspace{1pt}
\hspace{2pt}\textbf{Anglais} - Niveau B2


    \vspace{10pt}

    \subsection*{\textcolor{maincolor}{\faIcon{heart}} Intérêts}
\small

\hspace{4pt}\textbf{\color{maincolor}Tennis de table}
\begin{itemize}[leftmargin=11pt, itemsep=1pt, label={}]
    \item Compétition (niveau régional)
    \item Membre du comité directeur de mon club
\end{itemize}

\hspace{4pt}\textbf{\color{maincolor}Lecture}
\begin{itemize}[leftmargin=11pt, itemsep=1pt, label={}]
    \item Informatique, technique
    \item Dark fantaisie
    \item Manga
\end{itemize}

\hspace{4pt}\textbf{\color{maincolor}Jeux de société}
\begin{itemize}[leftmargin=11pt, itemsep=1pt, label={}]
    \item Jeux experts
    \item Campagne coopérative (ou non)
\end{itemize}


    % Les deux \vfill encadrant l'icône absorbent tout l'espace restant et la
    % centrent verticalement. Étant en glue infinie, ils sont seuls à s'étirer :
    % les espacements des listes ci-dessus gardent leur valeur naturelle.
    \vfill
    % ============================================================
% ICÔNE OMARCHY — SIDEBAR
% ============================================================
% Le dessin lui-même est dans config/icon.tex (partagé avec la variante une
% colonne, où il sert de filigrane centré sur chaque page).
%
% Placement : l'icône est simplement dans le flux, encadrée par les deux \vfill
% de sections/sidebar.tex, ce qui la centre dans l'espace libre sous le contenu
% de la sidebar. Volontairement PAS d'overlay « remember picture » : celui-ci
% demanderait deux passes de compilation, alors que le README documente un
% simple « pdflatex main.tex » — au premier passage l'icône serait mal placée.
%
% \omarchyiconsize : côté de l'icône (carrée). Borné par \linewidth, la largeur
% utile de la sidebar.
\newlength{\omarchyiconsize}
\setlength{\omarchyiconsize}{3.4cm}

% Teinte très diluée : filigrane à peine perceptible. On mélange la couleur au
% fond de la sidebar plutôt que d'utiliser « opacity », ce qui donne exactement
% le même rendu ici (le fond est un aplat uni) mais reste une couleur opaque
% ordinaire — pas de canal de transparence, donc aucun risque à l'impression ni
% dans les lecteurs PDF qui aplatissent mal les calques.
% \omarchyiconmix : part de maincolor, en pourcentage (0 = invisible).
\newcommand{\omarchyiconmix}{12}

{\centering
\omarchyicon{\omarchyiconsize}{maincolor!\omarchyiconmix!sidebarcolor}
\par}

    \vfill
    % \end{minipage} exécute un \unskip, qui supprimerait le \vfill ci-dessus
    % s'il terminait la boîte (l'icône se retrouverait collée en bas). Cette
    % \hbox vide, de dimensions nulles et sans interligne, le protège.
    \nointerlineskip\hbox{}%

\end{minipage}
